\documentclass[11pt]{article}
\usepackage[a4paper,margin=25mm]{geometry}
\usepackage{amsmath,amssymb,amsthm,mathtools,booktabs,microtype}
\usepackage[T1]{fontenc}\usepackage{lmodern}
\usepackage[colorlinks=true,linkcolor=blue,urlcolor=blue]{hyperref}
\newtheorem{theorem}{Theorem}[section]
\newtheorem{lemma}[theorem]{Lemma}
\newcommand{\F}{\mathbb F}
\newcommand{\SV}{\operatorname{SV}}
\title{Degree-Independent Direction Grids for\\
Width-Three Affine-Orbit Identity Testing}
\author{Yilin Liu \and GPT-6 Astra Pro}\date{26 October 2026}
\begin{document}\maketitle
\begin{abstract}
The nine-seed width-three affine-orbit construction admits direction
grids whose sizes do not depend on the latent individual degree.
The first preparation stage has at most \(n+1\) scalar conditions.
The final scalar decomposition has at most \(n+1\) factors.
Each relevant lowest homogeneous term has characteristic digit degree
at most two, giving a two-direction coefficient with at most \(n-1\)
distinct bad values in each sequentially chosen direction.
This yields \(O(n^9D^5)=O(n^{14}\delta^5)\) queries, where
\(D=n\delta\), over every characteristic. The structural nine-seed
theorem is a dependency; this note only refines its finite grid.
\end{abstract}

\noindent\textbf{Status.}
This quantitative argument has independent model-assisted internal
reconstructions and an exact-source audit. This is not external
peer review, and no publication-priority claim is made.

\section{Statement and structural dependency}

Let \(n,\delta\ge1\) and \(D=n\delta\).
The polynomial class consists of
\begin{equation}\label{eq:model}
 f(X)=u A_1(\ell_1(X)+c_1)\cdots A_n(\ell_n(X)+c_n)v,
\end{equation}
where the univariate matrices have width at most three, their entries
have degree at most \(\delta\), and the linear forms \(\ell_i\) are the
coordinates of an invertible linear transformation.
The representation, order, and affine transformation are unknown.

Fix \(n\) distinct abscissas in a working field \(\F\), and let
\(\lambda_i(y)\) be their Lagrange basis.
Set \(e_i=i\) in characteristic zero. In characteristic \(p>0\),
let \(e_i\) be the \(i\)-th positive integer not divisible by \(p\).
Write
\[
 \gamma_i(t)=t^{e_i},\qquad \rho=e_n\le2n-1,
 \qquad
 \mathcal H_i=\gamma_i(t)+\sum_{j=1}^4z_j\lambda_i(y_j).
\]

We use the independently reconstructed nine-seed theorem in
\emph{General Width-Three Affine Orbits via Mixed \(P_2/I_3\) Jets}
and its invertible-segment dependency.
The exact mathematical source reviewed for this deduction is
\path{width_three_general_affine_separating_curve_candidate.txt},
SHA-256 prefix \texttt{064b9e69}.
The properties needed here are recalled and counted below.
No interval-rank cone theorem is used.

\begin{theorem}\label{thm:main}
Choose arbitrary sets \(Y_1,\ldots,Y_4\subseteq\F\), each of size
\(n^2\), and sets
\[
 |T|=\rho D+1,\qquad |Z_j|=D+1\quad(1\le j\le4),
\]
with distinct elements within each set. The Cartesian evaluations
of \(\mathcal H\) on these sets hit every nonzero \eqref{eq:model}.
The number of oracle queries is at most
\begin{equation}\label{eq:bound}
 n^8(\rho D+1)(D+1)^4
 =O(n^9D^5)=O(n^{14}\delta^5).
\end{equation}
The field-size requirement
\[
 |\F|\ge\rho D+1
\]
suffices.
\end{theorem}

The sets need not be disjoint from one another or from the abscissas.
The evaluated map is polynomial, so every query is defined.
Over a smaller base field one may work in a sufficiently large
extension if oracle evaluations there are permitted.
If \(n=0\) or \(\delta=0\), one evaluation handles the constant case.

\section{A two-direction root bound}

In characteristic \(p>0\), the digit degree of a monomial is the sum
of the base-\(p\) digit sums of its exponents.
In characteristic zero use ordinary total degree.

\begin{lemma}\label{lem:witness}
Let \(K\supseteq\F\) be any coefficient-field extension.
Suppose a nonzero rational scalar \(g\), regular at a shifted origin,
has either a nonzero constant term or a lowest nonzero homogeneous
term \(H\) of digit degree at most two.
After two fresh SV pairs, there is a nonzero amplitude coefficient
\(C(y_1,y_2)\) of one of the following forms:
\begin{enumerate}
\item a nonzero constant;
\item a nonzero linear combination of \(\lambda_i(y_1)^a\);
\item
\[
 \sum_{i,j}c_{ij}\lambda_i(y_1)^a\lambda_j(y_2)^b\ne0.
\]
\end{enumerate}
Here each positive exponent \(a,b\) is one in characteristic zero
or a pure \(p\)-power in characteristic \(p\).
The claim remains true after an invertible linear change of the
residual variables.
\end{lemma}
\begin{proof}
Linear substitution preserves the digit-degree upper bound and
preserves nonzeroness when invertible.
If the constant term is nonzero, take that coefficient.
Otherwise \(H\) is homogeneous of positive total degree.
A pure-power sum gives case2.
For degree \(p^e+p^f\) with \(e\ne f\), the coefficient of
\(z_1^{p^e}z_2^{p^f}\) gives case3 and is nonzero by independence
of the two Frobenius-powered Lagrange bases.
For degree \(2p^e\) in odd characteristic, quadratic polarization
gives a nonzero coefficient of this form.
In characteristic two, nonzero cross terms give the nonzero
off-diagonal bilinear coefficient; if there are no cross terms,
the pure-power case applies.
The characteristic-zero cases are linear and quadratic.
These exhaust the possibilities. Higher homogeneous terms have
larger total amplitude degree and cannot affect the chosen coefficient.
\end{proof}

\begin{lemma}[Simultaneous sequential directions]\label{lem:simultaneous}
Suppose \(M\) nonzero coefficients \(C_\nu(y_1,y_2)\) have the forms
in Lemma~\ref{lem:witness}, with exponents permitted to differ
between coefficients. Any two sets of at least \(M(n-1)+1\)
distinct elements of \(\F\) contain a pair \((\eta_1,\eta_2)\)
such that every \(C_\nu(\eta_1,\eta_2)\) is nonzero in \(K\).
\end{lemma}
\begin{proof}
For a nonzero case3 coefficient, expand in the independent basis
\(\{\lambda_j(y_2)^b\}\). At least one coefficient is a nonzero
combination of \(\lambda_i(y_1)^a\).
Over an algebraic closure of \(K\), this combination is an \(a\)-th
power of a nonzero polynomial of degree at most \(n-1\).
It therefore has at most \(n-1\) distinct roots.
Thus at most \(n-1\) choices of \(\eta_1\) make the entire
coefficient vanish identically in \(y_2\).
Case2 has the same root bound; a constant has no bad choices.

Choose \(\eta_1\) outside the union of all these bad sets.
For every remaining nonconstant coefficient, its specialization is
a nonzero combination of \(\lambda_j(y_2)^b\), so the same argument
excludes at most \(n-1\) second directions. Choose \(\eta_2\)
outside their union.
Root extraction is used only in the algebraic-closure proof of
the distinct-root bound. No perfectness or effective root extraction
is required.
\end{proof}

If the rational scalar has denominators that are units at the shifted
origin, these remain units after every fixed direction choice: their
constant amplitude terms do not depend on the directions.
Thus a nonzero coefficient supplied by the lemmas certifies a
well-defined nonzero rational substitution.

\section{Counting the first-stage scalar conditions}

Assume \(f\ne0\), so no original local factor is zero.
Let \(r\) be the number of original factors of rank one and \(q\)
the number of rank two. In particular \(r+q\le n\).
First split the word at its rank-one factors by nonzero univariate
pivots. This gives \(r+1\) nonzero scalar chains.
Let \(c\) chains contain at least one rank-two factor and
\(c_0=r+1-c\) contain none.

For a chain with \(k\ge1\) rank-two factors, choose rational
univariate full-rank factorizations \(A_i=U_iV_i\).
Writing \(M_i\) for its intervening rank-three runs, define
\[
 \alpha=\mathbf lM_0U_1,\qquad
 \beta=V_kM_k\mathbf r,\qquad
 T_i=V_iM_iU_{i+1}\quad(1\le i<k).
\]
Here \(\mathbf l,\mathbf r\) are the chain's endpoint vectors.
It suffices to preserve the following chosen nonzero scalar functions:
\[
\begin{array}{c|c}
\text{condition in one chain}&\text{maximum count}\\ \hline
\text{one entry of }\alpha\text{ and one of }\beta&2\\
\text{one pivot for each rank-one }T_i&
 \#\{i:\operatorname{rank}T_i=1\}\\
\det T_i\text{ for each rank-two }T_i&
 \#\{i:\operatorname{rank}T_i=2\}
\end{array}
\]
Every \(T_i\) has rank one or two: a zero transfer would make the
original scalar chain zero. Thus each transfer contributes exactly
one test, and the total in this chain is \(k+1\).
For a chain with no rank-two factor, preserve its single nonzero
scalar output.
Consequently the total number of tests satisfies
\begin{equation}\label{eq:counttests}
 M\le q+c+c_0=q+r+1\le n+1.
\end{equation}

A chosen entry keeps each endpoint vector nonzero.
A chosen rank-one pivot keeps its designated row and column nonzero,
and its transfer remains rank one because all two-by-two minors
vanish identically. A chosen rank-two determinant keeps its transfer
invertible. These are exactly the facts used by the scalar splitting
and the active-boundary gauges; no other transfer entry is required
to stay nonzero.

Every entry test is an invertible three-state run with at most two
rational univariate boundary vectors on distinct variables.
A determinant test has the same form after taking second compounds:
\[
 \det(V_iM_iU_{i+1})
   =(\wedge^2V_i)(\wedge^2M_i)(\wedge^2U_{i+1}).
\]
The middle compound matrices are still invertible three-by-three
matrices. The restricted-rank theorem gives a lowest homogeneous
term of digit degree at most two for each test that vanishes on
the initial curve. Tests nonzero on that curve need no bad-direction
exclusion.

There are no omitted scalar conditions for local pivots, endpoint
completions, or exterior-square denominators.
Their denominators and nonzero determinants are univariate rational
functions chosen before any seed substitution.
For \emph{every} fixed pair of direction values, and with their
amplitudes still symbolic,
\[
 C_i(t)=c_i+\ell_i\bigl(\gamma(t)+
 z_1\lambda(\eta_1)+z_2\lambda(\eta_2)\bigr)
\]
has derivative \(\ell_i(\gamma'(t))\ne0\).
Thus \(C_i\) is transcendental over \(\F(z_1,z_2)\), so every
nonzero such univariate rational function remains defined and
nonzero. This covers the original rank-one pivots, the two-by-two
pivots defining \(U_i,V_i\), their compound denominators, and the
standard-basis endpoint completions.

Apply Lemmas~\ref{lem:witness}--\ref{lem:simultaneous} over \(K=\F(t)\)
to the at most \(n+1\) scalar tests.
Since \((n+1)(n-1)+1=n^2\), the first two grids contain old directions for
which all tests survive as rational functions in \(t,z_1,z_2\).
All transfers then have their generic ranks, and both boundary
active projections are nonzero.

\section{Counting the final scalar factors}

Fix the old directions just obtained and keep \(z_1,z_2\) symbolic.
Split each rank-one transfer \(T_i\) at a nonzero entry.
Use the same chosen pivots as in the preceding tests, so their values
are nonzero at the new shifted origin.
There are at most \(q-c\) internal transfers across the \(c\)
chains with rank-two factors. If \(m_1\) of them have rank one,
then \(m_1\le q-c\).
The final number of nonzero rational scalar factors is therefore
\begin{equation}\label{eq:factorcount}
 (r+1)+m_1\le r+q-c+1\le n+1.
\end{equation}

The structural theorem recombines every adjacent \(U_iV_i\) into
the original univariate factor. Thus each scalar numerator still
has one occurrence of each of its latent variables, even though
different scalar factors may share variables.
Every mixed segment has exact constant matrices \(P_2\) or \(I_3\)
after cut gauges, with nonzero active boundary projections.
The theorem proves that its entire lowest homogeneous term has
digit degree at most two. Division by a transfer pivot only scales
that term by the nonzero inverse constant of the pivot.
A factor already nonzero at the shifted origin has no bad directions.

Apply Lemma~\ref{lem:simultaneous} over
\(K=\F(z_1,z_2)(t)\) to these at most \(n+1\) factors.
Since
\[
 (n+1)(n-1)+1=n^2,
\]
the two new direction grids contain a pair preserving every factor.
Transfer-pivot denominators have nonzero constants in the new
amplitudes and need no additional exclusion. Cut gauges are
constructed after the directions are fixed; their nonzero minors
follow from the preserved rank and boundary conditions, rather than
being extra tests.

The product is consequently nonzero as a rational function in
\(t,z_1,z_2,z_3,z_4\). The exact factorization identity implies that
the original polynomial \(f(\mathcal H)\), with all four directions
fixed, is a nonzero polynomial in those five variables.

\section{Finish the grid proof}

\begin{proof}[Proof of Theorem~\ref{thm:main}]
The two preceding sections choose, inside the prescribed direction
grids, a tuple for which the original polynomial remains nonzero.
After that tuple is fixed, the degree in \(t\) is at most \(\rho D\),
and the degree in each amplitude is at most \(D\).
The usual individual-degree interpolation bound gives a nonzero
evaluation on \(T\times Z_1\times\cdots\times Z_4\).
The full Cartesian search includes these adaptively established
direction choices and therefore hits every input polynomial.

This proof uses rational functions only to certify symbolic
nonzeroness. It imposes no denominator restriction on the final
evaluation points: the oracle is queried on the original polynomial.
The number of queries is \eqref{eq:bound}.
Finally \(\rho\ge n\) and \(D=n\delta\ge n\), so
\(\rho D+1\ge n^2+1\), as well as \(D+1\) and \(n\).
The field-size assumption supplies all required sets and abscissas.
\end{proof}

One may instead use the ordinary individual-degree grid for the
same nine-seed map. The present theorem is a separately justified
choice of grids; it does not assert that arbitrary coordinatewise
mixtures of these two grid prescriptions inherit both proofs.
\end{document}
